% Template KLTN cho SV trường ĐHKHTN
% Liên hệ: nqminh@fit.hcmus.edu.vn
% Last update: 08/06/2016

% Chú ý: đọc các phần chú ý đóng khung của file này và chỉnh lại cho phù hợp.
% Trước khi build, xóa hết các file được tạo ra trong quá trình build trước đó, và build theo thứ tự: BIB > PDF > PDF.
% Nếu cập nhật tài liệu tham khảo, cũng cần build lại theo cách trên.

\documentclass[oneside,a4paper,14pt]{extreport}

% Font tiếng Việt
\usepackage[T5]{fontenc}
\usepackage[utf8]{inputenc}
\DeclareTextSymbolDefault{\DH}{T1}

% Tài liệu tham khảo
\usepackage[
	sorting=nty,
	backend=bibtex,
	defernumbers=true]{biblatex}
\usepackage[unicode]{hyperref} % Bookmark tiếng Việt
\addbibresource{References/references.bib}

\makeatletter
\def\blx@maxline{77}
\makeatother

% Chèn hình, các hình trong luận văn được để trong thư mục Images/
\usepackage{graphicx}
\graphicspath{ {Images/} }

% Chèn và định dạng mã nguồn
\usepackage{listings}
\usepackage{color}
\definecolor{codegreen}{rgb}{0,0.6,0}
\definecolor{codegray}{rgb}{0.5,0.5,0.5}
\definecolor{codepurple}{rgb}{0.58,0,0.82}
\definecolor{backcolour}{rgb}{0.95,0.95,0.92}
\lstdefinestyle{mystyle}{
    backgroundcolor=\color{backcolour},   
    commentstyle=\color{codegreen},
    keywordstyle=\color{magenta},
    numberstyle=\tiny\color{codegray},
    stringstyle=\color{codepurple},
    basicstyle=\footnotesize,
    breakatwhitespace=false,         
    breaklines=true,                 
    captionpos=b,                    
    keepspaces=true,                 
    numbers=left,                    
    numbersep=5pt,                  
    showspaces=false,                
    showstringspaces=false,
    showtabs=false,                  
    tabsize=2
}
\lstset{style=mystyle}

% Chèn và định dạng mã giả
\usepackage{amsmath}
\usepackage{algorithm}
\usepackage[noend]{algpseudocode}
\makeatletter
\def\BState{\State\hskip-\ALG@thistlm}
\makeatother

% Bảng biểu
\usepackage{multirow}
\usepackage{array}
\newcolumntype{L}[1]{>{\raggedright\let\newline\\\arraybackslash\hspace{0pt}}m{#1}}
\newcolumntype{C}[1]{>{\centering\let\newline\\\arraybackslash\hspace{0pt}}m{#1}}
\newcolumntype{R}[1]{>{\raggedleft\let\newline\\\arraybackslash\hspace{0pt}}m{#1}}



 
% Đổi tên mặc định
\renewcommand{\chaptername}{Chương}
\renewcommand{\figurename}{Hình}
\renewcommand{\tablename}{Bảng}
\renewcommand{\contentsname}{Mục lục}
\renewcommand{\listfigurename}{Danh sách hình}
\renewcommand{\listtablename}{Danh sách bảng}
\renewcommand{\appendixname}{Phụ lục}

% Định dạng chapter
\usepackage{titlesec}
\titleformat{\chapter}
    [display] 
    {\normalfont\bfseries\Large}{\chaptername \ \thechapter}{10pt}{\huge}
\titlespacing*{\chapter}{0pt}{-10pt}{40pt} %khoảng cách giữa chapter và đầu trang

\titleformat{\section}
    {\normalfont\bfseries\large}{\thesection}{1em}{}
 
\titleformat{\subsection}
    {\normalfont\bfseries\normalsize}{\thesubsection}{1em}{}
  
% Dãn dòng 1.5
\usepackage{setspace}
\onehalfspacing

% Thụt vào đầu dòng
\usepackage{indentfirst}

% Canh lề
\usepackage[
  top=20mm,
  bottom=10mm,
  left=30mm,
  right=20mm,
  footskip = 15mm,
  includefoot]{geometry}
  
% Trang bìa
\usepackage{tikz}
\usetikzlibrary{calc}
\newcommand\HRule{\rule{\textwidth}{1pt}}

% ========================================================================================= %
% CHÚ Ý: Thông tin chung về KLTN - sinh viên điền vào đây để tự động update các trang khác  %
% ========================================================================================= %
\newcommand{\tenSV}{Nguyễn~Văn~A~-~Trần~Văn~B} % Dấu ~ là khoảng trắng không được tách (các chữ nối với nhau bằng dấu ~ sẽ nằm cùng 1 dòng
\newcommand{\mssv}{1234567}
\newcommand{\tenKL}{Sử~dụng~LaTeX trong Khoá~luận~tốt~nghiệp} % Chú ý dấu ~ trong tên khóa luận
\newcommand{\tenGVHD}{Tên~Giáo~Viên}
\newcommand{\tenBM}{Công nghệ tri thức}

\begin{document}

\begin{titlepage}

    \begin{center}
    %ĐẠI HỌC QUỐC GIA THÀNH PHỐ HỒ CHÍ MINH\\
    TRƯỜNG ĐẠI HỌC KHOA HỌC TỰ NHIÊN\\
    \textbf{KHOA CÔNG NGHỆ THÔNG TIN}\\[2cm]


    { \Large \bfseries \tenSV\\[2cm] } 

    %Tên đề tài Khóa luận tốt nghiệp/Đồ án tốt nghiệp

    { \Large \bfseries Meta-Learning và Personalization Layer trong Federated Learning\\[3cm]} 


    %Chọn trong các dòng sau
    \large KHÓA LUẬN TỐT NGHIỆP CỬ NHÂN\\
    %\large ĐỒ ÁN TỐT NGHIỆP CỬ NHÂN\\
    %\large THỰC TẬP TỐT NGHIỆP CỬ NHÂN\\
    %Đưa vào dòng này nếu thuộc chương trình Chất lượng cao, hoặc lớp Cử nhân tài năng
    \large CHƯƠNG TRÌNH CHÍNH QUY\\
    %\large CHƯƠNG TRÌNH CHẤT LƯỢNG CAO\\
    %\large CHƯƠNG TRÌNH CỬ NHÂN TÀI NĂNG\\[2cm]


    \begin{tikzpicture}[remember picture, overlay]
        \draw[line width = 2pt] ($(current page.north west) + (2cm,-2cm)$) rectangle ($(current page.south east) + (-1.5cm,2cm)$);
    \end{tikzpicture}

    \vfill
    Tp. Hồ Chí Minh, tháng 02/2022

    \end{center}

    \pagebreak

    \begin{center}

    TRƯỜNG ĐẠI HỌC KHOA HỌC TỰ NHIÊN\\
    \textbf{KHOA CÔNG NGHỆ THÔNG TIN}\\[2cm]

    {\large \bfseries Nguyễn Bảo Long - 18120201\\} 
    {\large \bfseries Cao Tất Cường - 18120296\\[2cm]}

    %Tên đề tài Khóa luận tốt nghiệp/Đồ án tốt nghiệp
    { \Large \bfseries Meta-Learning và Personalization Layer trong Federated Learning\\[3cm]} 


    %Chọn trong các dòng sau
    \large KHÓA LUẬN TỐT NGHIỆP CỬ NHÂN\\
    %\large ĐỒ ÁN TỐT NGHIỆP CỬ NHÂN\\
    %Đưa vào dòng này nếu thuộc chương trình Chất lượng cao, hoặc lớp Cử nhân tài năng
    \large CHƯƠNG TRÌNH CHÍNH QUY\\[2cm]
    %\large CHƯƠNG TRÌNH CHẤT LƯỢNG CAO\\[2cm]
    %\large CHƯƠNG TRÌNH CỬ NHÂN TÀI NĂNG\\[2cm]

    \textbf{GIẢNG VIÊN HƯỚNG DẪN}\\
    \tenGVHD

    \begin{tikzpicture}[remember picture, overlay]
        \draw[line width = 2pt] ($(current page.north west) + (2cm,-2cm)$) rectangle ($(current page.south east) + (-1.5cm,2cm)$);
    \end{tikzpicture}

    \vfill
    Tp. Hồ Chí Minh, tháng 02/2022

    \end{center}

\end{titlepage}
% Sasu trang Title, các bạn chèn nhận xét gủa GVHD và GVPB. Nhận xét sẽ được giáo vụ phát sau buổi bảo vệ để các bạn đóng quyển.

\pagenumbering{roman} % Đánh số i, ii, iii, ...

%\addcontentsline{toc}{chapter}{Lời cam đoan}
%\chapter*{Lời cam đoan}
\label{reassurances}

Tôi xin cam đoan đây là công trình nghiên cứu của riêng tôi. Các số liệu và kết quả nghiên cứu trong luận văn này là trung thực và không trùng lặp với các đề tài khác.

\addcontentsline{toc}{chapter}{Lời cảm ơn}
\chapter*{Lời cảm ơn}
\label{thanks}

Tôi xin chân thành cảm ơn ...

\addcontentsline{toc}{chapter}{Đề cương chi tiết}
\include{Appendix/decuong}

% Mục lục, danh sách hình, danh sách bảng
\addcontentsline{toc}{chapter}{Mục lục}
\tableofcontents
\listoffigures
\listoftables

\addcontentsline{toc}{chapter}{Tóm tắt}
\include{Appendix/tomtat}

\clearpage

\pagenumbering{arabic} % Đánh số 1, 2, 3, ...

% Các chương nội dung
\chapter{Giới thiệu}
\label{Chapter1}

\section{Đặt vấn đề \& Động lực}

Hiện nay, các thiết bị biên như điện thoại, máy tính bảng, thậm chí máy giặt, máy hút bụt thông minh có thể sinh ra lượng lớn dữ liệu trong quá trình hoạt động. Lượng dữ liệu này, nếu tận dụng được, có thể mang lại sự cải thiện rất lớn về độ chính xác cho các mô hình máy học hiện tại. Ví dụ, dữ liệu thu thập được từ bàn phím điện thoại có thể phục vụ tối ưu cho các mô hình ngôn ngữ; ảnh chụp được lưu trữ trong bộ nhớ điện thoại hoàn toàn có thể được sử dụng làm dữ liệu để huấn luyện cho mô hình nhận dạng ảnh; hay lịch sử duyệt web của người dùng có thể được dùng cho bài toán đề xuất sản phẩm. Những lý do trên trở thành một động lực to lớn, thúc đẩy việc tìm ra một phương pháp giúp tận dụng nguồn dữ liệu dồi dào này.

Việc ngày càng nhiều dữ liệu được sinh ra tại các thiết bị biên khiến cho phương pháp huấn luyện mô hình theo cách tiếp cận truyền thống (được gọi là huấn luyện tập trung) bộc lộ nhiều khuyết điểm. Ba điểm yếu khiến cho cách tiếp cận này không còn mạnh mẽ có thể kể đến: (1) - Sự vi phạm về quyền riêng tư dữ liệu, (2) - Chi phí truyền tin, (3) - Chi phí phần cứng máy chủ.

\textbf{Sự vi phạm về quyền riêng tư dữ liệu.} Phương pháp truyền thống đòi hỏi phải gửi dữ liệu người dùng về một máy chủ để tiến hành huấn luyện mô hình. Các thông tin nhạy cảm của người dùng hoàn toàn có thể bị nghe lén bởi kẻ tấn công hoặc bị khai thác khi máy chủ bị nhiễm mã độc. Điều này ảnh hưởng nghiêm trọng đến quyền riêng tư dữ liệu của người dùng - một vấn đề mà hiện nay đang nhận được rất nhiều sự quan tâm từ cả người dùng lẫn chính phủ.

\textbf{Chi phí truyền tin.} Một người dùng điện thoại thông minh giờ đây có thể thực hiện giao dịch tài chính, lướt web, xem phim ngay trên thiết bị của mình khiến cho dữ liệu được sinh ra liên tục. Một máy hút bụi thông minh được trang bị các cảm biến nên dữ liệu cũng được sinh ra liên tục trong quá trình vận hành. Điều này khiến cho dữ liệu tại thiết bị biên ngày một tăng lên. Chi phí truyền tin từ các thiết bị biên đến máy chủ để huấn luyện trở nên tốn kém và có thể gây mất thông tin, ảnh hưởng đến hiệu suất học của mô hình.

\textbf{Chi phí phần cứng máy chủ.} Sau khi dữ liệu được gửi về máy chủ, cần một cấu hình máy mạnh mẽ cùng khả năng lưu trữ lớn để có thể xử lý hết lượng dữ liệu khổng lồ trên trong thời gian giới hạn.

Việc các phương pháp tiếp cận máy học theo hướng truyền thống đang dần bộc lộ các nhược điểm về chi phí vận hành và bảo trì ngày càng cao, cũng như các mối nguy hiểm tiềm tàng có thể xảy ra đối với dữ liệu của người dùng, một lần nữa thúc đẩy việc nghiên cứu về một phương pháp huấn luyện giúp làm giảm chi phí phần cứng (sử dụng cho đường truyền và máy chủ), đồng thời đảm bảo tính riêng tư dữ liệu cho người dùng.

Khái niệm \textit{federated learning} và thuật toán \codeword{FedAvg} được đưa ra vào năm 2016 bởi Google trong nghiên cứu \cite{mcmahan2017communication} nhằm mục đích huấn luyện mô hình máy học trên các tập dữ liệu riêng biệt được phân bố trên các thiết bị biên (được gọi là huấn luyện phân tán). Do đó, \textbf{một hệ thống FL không cần một máy chủ quá mạnh để vận hành (thậm chí có thể sử dụng một máy khách để vận hành \cite{yin2021comprehensive}), không đòi hỏi chi phí truyền tin quá lớn và đảm bảo được quyền riêng tư dữ liệu của người dùng vì không diễn ra bất cứ quá trình thu thập dữ liệu từ người dùng nào (điều mà mô hình huấn luyện tập trung bắt buộc phải làm)}. Dễ thấy rằng, phần lớn quá trình tính toán được phân tán đến các thiết bị biên. Tuy nhiên, khả năng lưu trữ và tính toán tại các thiết bị này ngày càng được cải thiện, khiến cho việc huấn luyện phân tán dần trở nên khả thi và đạt hiệu quả cao hơn. Bằng chứng là sự xuất hiện của hệ thống FL trong nhiều lĩnh vực như Internet vạn vật (IoT), y tế và sức khỏe, các hệ thống đề xuất, chuỗi cung ứng,... \cite{zhou2021survey}

Bên cạnh đó, nghiên cứu \cite{zhao2018federated} chỉ ra rằng, hệ thống FL hoạt động trên nền thuật toán \codeword{FedAvg} bị giảm hiệu suất nghiêm trọng khi xử lý dữ liệu Non-IID. Tuy nhiên, dữ liệu phân bố trên máy khách là không đồng nhất và có tính cá nhân hóa tất cao. Nói cách khác, các tập dữ liệu này tuân theo phân phối Non-IID. Do đó, \textbf{khóa luận này được thực hiện nhằm nghiên cứu và cải tiến hệ thống FL để thu được kết quả cao cũng như cải thiện khả năng cá nhân hóa mô hình học cho từng người dùng trên dữ liệu Non-IID}.

\section{Phạm vi đề tài}

Nghiên cứu \cite{yin2021comprehensive} chỉ ra ba hướng nghiên cứu chính khi đề cập đến một hệ thống FL: (1) - Cải thiện hiệu suất của hệ thống FL, (2) - Cải thiện khả năng bảo mật của hệ thống FL, (3) - Cải thiện vấn đề về quyền riêng tư của người dùng trong hệ thống FL.

Về việc phân loại hệ thống FL, dựa trên dữ liệu đầu vào, hệ thống FL được chia thành ba loại \cite{yin2021comprehensive}: (1) - Horizontal FL, (2) - Vertical FL, (3) - Federated transfer learning.

Về việc phân loại các kịch bản Non-IID, nghiên cứu \cite{zhu2021federated} chỉ ra bốn kịch bản chính: (1) - Phân phối thuộc tính khác nhau giữa các máy khách, (2) - Phân phối nhãn khác nhau giữa các máy khách, (3) - Phân phối thời gian khác nhau giữa các máy khách, (4) - Các kịch bản khác.

Từ đó, khóa luận này có phạm vi nghiên cứu được giới hạn và phương án giải quyết được xây dựng dựa trên ba giả định sau:

\begin{itemize}
    \item Loại hệ thống: Môi trường thí nghiệm (bao gồm các yếu tố như số lượng người dùng, dữ liệu, cấu hình, khả năng lưu trữ của thiết bị cuối,...) tuân theo đặc trưng của hệ thống Horizontal FL.
    \item Bảo mật \& quyền riêng tư: Hệ thống đã đảm bảo tính bảo mật cũng như duy trì tốt quyền riêng tư của người dùng.
    \item Kịch bản Non-IID: Phân phối nhãn dữ liệu trên các máy khách là khác nhau.
\end{itemize}

\section{Đóng góp chính}

Các đóng góp chính của khóa luận được chia thành hai loại: đóng góp về mặt lý thuyết và đóng góp về mặt thực nghiệm.

\subsection{Đóng góp lý thuyết}

\begin{itemize}
    \item Nghiên cứu hệ thống FL và thách thức về phân phối dữ liệu mà hệ thống Horizontal FL gặp phải.
    \item Khảo sát các phương pháp tối ưu hóa hệ thống Horizontal FL trên dữ liệu dạng Non-IID. Trong đó, tập trung nghiên cứu các phương pháp theo hướng Personalized Federated Averaging \parencite{fallah2020personalized, chen2018federated} và Personalization Layer \parencite{liang2020think, arivazhagan2019federated}.
    \item Phương pháp đề xuất của khóa luận đã cho thấy khả năng đạt được độ chính xác cao hơn trong quá trình kiểm thử với hai đối tượng người dùng (người dùng cục bộ và người dùng mới) so với các phương pháp trước đó (chỉ sử dụng \codeword{FedAvg}, chỉ sử dụng Personalized Federated Averaging, hoặc chỉ sử dụng Personalization Layer).
\end{itemize}

\subsection{Đóng góp thực nghiệm}

\begin{itemize}
    \item Tổ chức bộ dữ liệu MNIST và CIFAR-10 theo hai hướng IID và Non-IID để tiến hành thí nghiệm.
    \item Cài đặt thuật toán \codeword{FedAvg}, \codeword{FedAvgMeta}, các thuật toán kết hợp giữa \codeword{FedAvg} và ML (thuật toán \codeword{FedMeta(MAML)}, \codeword{FedMeta(Meta-SGD)}).
    \item Cài đặt thuật toán kết hợp giữa \codeword{FedAvg} và PL (thuật toán \codeword{FedPer}, \codeword{LG-FedAvg}).
    \item Kết hợp các thuật toán ML và PL vào hệ thống FL: các thuật toán \codeword{FedMeta-Per}.
    \item Fine-tune các siêu tham số như số lượng máy khách tham gia huấn luyện, số bước huấn luyện cục bộ, các siêu tham số học để mô hình đạt độ chính xác tốt nhất.
\end{itemize}

\section{Bố cục}

Trong luận văn này, chương \ref{Chapter2} trình bày về tổng quan lý thuyết được sử dụng trong khóa luận, các lý thuyết này làm nền tảng cho nghiên cứu và đề xuất thuật toán; chương \ref{Chapter3} đề xuất thuật toán giúp giải quyết vấn đề vừa nêu ở chương \ref{Chapter1}; chương \ref{Chapter4} trình bày về cài đặt thực nghiệm để kiểm chứng tính hiệu quả của thuật toán; chương \ref{Chapter5} đi vào phân tích kết quả đạt được; chương \ref{Chapter6} nêu kết luận và hướng phát triển tương lai của khóa luận.


\chapter{Trình bày báo cáo}
\label{Chapter2}

\section{Quy định chung}

Báo cáo phải được trình bày ngắn gọn, rõ ràng, mạch lạc, sạch sẽ, không được tẩy xóa, có đánh số trang, đánh số bảng biểu, hình vẽ, đồ thị. 

Nội dung báo cáo được phân thành các chương. Số thứ tự của các chương, mục được đánh số bằng hệ thống số Ả-rập, không dùng số La mã. Các mục và tiểu mục được đánh số bằng các nhóm hai hoặc ba chữ số, cách nhau một dấu chấm: số thứ nhất chỉ số chương, chỉ số thứ hai chỉ số mục, số thứ ba chỉ số tiểu mục.


%Báo cáo cần dùng LaTEX để viết và trình bày theo mẫu đã được cung cấp.

 Báo cáo trình bày sử dụng khổ giấy với việc canh lề như sau: Lề trên 3 cm, lề dưới 2,5 cm, lề trái 3 cm, lề phải 2 cm. Đánh số trang ở giữa bên dưới. Đánh số trang ở giữa bên dưới.

Font chữ dùng trong báo cáo (Times New Roman) với kích cỡ (size) 13-14pt, sử dụng chế độ dãn dòng (line spacing) chế độ 1.5 lines.

%Các bảng biểu trình bày theo chiều ngang khổ giấy thì đầu bảng là lề trái của trang. 


\section{Bố cục của báo cáo}

Nội dung của báo cáo tối thiểu 50 trang khổ A4 và không nên vượt quá 100 trang (không kể các trang bìa, lời cám ơn, mục lục, tài liệu tham khảo \ldots) theo trình tự như sau:

\begin{itemize}
\item MỞ ĐẦU (thường đặt tên là ``Giới thiệu''): Trình bày lý do chọn đề tài, mục đích, đối tượng và phạm vi nghiên cứu.
Mô tả bài toán mà đề tài giải quyết.
Bài toán này có gì hay?
Tại sao lại cần giải quyết bài toán này?
Bài toán này có gì khó?
Có những hướng nào để giải quyết bài toán này?
Những hướng giải quyết trước đây có những vấn đề gì chưa giải quyết được?
Các câu hỏi nghiên cứu mà đề tài trả lời hoặc những vấn đề mà đề tài sẽ giải quyết.
Các đóng góp của đề tài.

\item TỔNG QUAN (thường đặt tên là ``Các công trình liên quan''): Phân tích đánh giá các hướng nghiên cứu đã có của các tác giả trong và ngoài nước liên quan đến đề tài; nêu những vấn đề còn tồn tại (những vấn đề nào mà các công trình khác chưa giải quyết được); chỉ ra những vấn đề mà đề tài cần tập trung, nghiên cứu giải quyết.

\item NGHIÊN CỨU THỰC NGHIỆM HOẶC LÝ THUYẾT (thường đặt tên là ``Phương pháp đề xuất''): Trình bày cơ sở lý thuyết, lý luận, giả thiết khoa học và phương pháp nghiên cứu đã được sử dụng trong đề tài.

Nếu đề xuất hướng giải quyết mới, mô hình mới thì cần mô tả chi tiết cách giải quyết của mình (chi tiết tới mức người khác có thể dựa vào phần này mà cài đặt lại được đúng hoàn toàn phương pháp của mình đề ra).

\item TRÌNH BÀY, ĐÁNH GIÁ BÀN LUẬN VỀ CÁC KẾT QUẢ (thường đặt tên là ``Kết quả thí nghiệm''): Mô tả các kết quả nghiên cứu khoa học hoặc kết quả thực nghiệm.
Đối với  đề tài ứng dụng có kết quả là sản phẩm phần mềm phải có hồ sơ thiết kế, cài đặt,\ldots theo một trong các mô hình đã học (UML,\ldots).

Thông thường cần mô tả môi trường thí nghiệm trước như sử dụng dữ liệu nào, dùng độ đo nào để đánh giá, môi trường chạy thí nghiệm (cấu hình máy nếu cần phân tích thông tin về thời gian chạy thực nghiệm). Sau đó, nêu kết quả thực nghiệm, bàn luận và giải thích kết quả.

\item KẾT LUẬN VÀ HƯỚNG PHÁT TRIỂN (thường đặt tên là ``Kết luận''): Trình bày những kết quả đạt được, những đóng góp mới và những đề xuất mới, kiến nghị về những hướng nghiên cứu tiếp theo.

\item DANH MỤC TÀI LIỆU THAM KHẢO: Chỉ bao gồm các tài liệu được trích dẫn, sử dụng và đề cập tới để bàn luận trong báo cáo.
Phần này các bạn chuẩn bị 1 file BIB để lưu các tài liệu trích dẫn.
Khi các bạn trích dẫn một tài liệu nào đó, LaTeX sẽ tự động thêm vào danh mục tài liệu tham khảo giúp các bạn.
Các bạn xem hướng dẫn cách trích dẫn ở chương sau.

\item PHỤ LỤC: Phần này bao gồm nội dung cần thiết nhằm minh họa hoặc hỗ trợ cho nội dung báo cáo như số liệu, mẫu biểu, tranh ảnh,\ldots Phụ lục không được dày hơn phần chính của báo cáo.
Nếu có công trình công bố thì để vào phần phụ lục này.
\end{itemize}

\section{Bảng biểu, hình vẽ, phương trình}

%Những qui định dưới này các bạn có thể bỏ qua hoặc đọc để hiểu thêm.
%Những định dạng này LaTeX đều tự động giúp các bạn.
%Các bạn xem hướng dẫn chi tiết hơn ở chương sau.

Việc đánh số bảng biểu, hình vẽ, phương trình phải gắn với số chương; ví dụ hình 3.4 có nghĩa là hình thứ 4 trong Chương 3.
Mọi đồ thị, bảng biểu, hình vẽ lấy từ các nguồn khác phải được trích dẫn đầy đủ.

\subsection{Bảng biểu, hình vẽ}



Đầu đề của bảng biểu ghi phía trên bảng, đầu đề của hình vẽ ghi phía dưới hình.

Thông thường, những bảng ngắn và đồ thị phải đi liền với phần nội dung đề cập tới các bảng và đồ thị này ở lần thứ nhất.
Các bảng dài có thể để ở những trang riêng nhưng cũng phải tiếp theo ngay phần nội dung đề cập tới bảng này ở lần đầu tiên.
Các bảng rộng vẫn nên trình bày theo chiều đứng dài 297mm của trang giấy, chiều rộng của trang giấy có thể hơn 210mm.
Chú ý gấp trang giấy sao cho số và đầu đề của hình vẽ hoặc bảng vẫn có thể nhìn thấy ngay mà không cần mở rộng tờ giấy.
Tuy nhiên hạn chế sử dụng các bảng quá rộng này.

Đối với những trang giấy có chiều đứng hơn 297mm (bản đồ, bản vẽ,\ldots) thì có thể để trong một phong bì cứng đính bên trong bìa sau của báo cáo.

Các hình vẽ phải sạch sẽ bằng mực đen để có thể sao chụp lại; có đánh số và ghi đầy đủ đầu đề, cỡ chữ phải bằng cỡ chữ sử dụng trong báo cáo.

Khi đề cập đến các bảng biểu và hình vẽ phải nêu rõ số của hình và bảng biểu đó, ví dụ ``... được nêu trong Bảng 4.1'' hoặc ``xem Hình 3.2'' mà không được viết ``… được nêu trong bảng dưới đây'' hoặc ``trong đồ thị của X và Y sau''.

\subsection{Phương trình toán học}

Việc trình bày phương trình toán học trên một dòng đơn hoặc dòng kép tùy ý, tuy nhiên phải thống nhất trong toàn báo cáo.

Khi ký hiệu xuất hiện lần đầu tiên thì phải giải thích và đơn vị tính phải đi kèm ngay trong phương trình có ký hiệu đó.
Nếu cần thiết, danh mục của tất cả các ký hiệu, chữ viết tắt và nghĩa của chúng cần được liệt kê và để ở phần đầu của báo cáo.

Tất cả các phương trình cần được đánh số và để trong ngoặc đơn đặt bên phía lề phải.
Nếu một nhóm phương trình mang cùng một số thì những số này cũng được để trong ngoặc, hoặc mỗi phương trình trong nhóm phương trình (5.1) có thể được đánh số là (5.1.1), (5.1.2), (5.1.3).

\section{Viết tắt}

\textbf{Không lạm dụng việc viết tắt} trong báo cáo.
Chỉ viết tắt những từ, cụm từ hoặc thuật ngữ được sử dụng nhiều lần trong báo cáo.
Không viết tắt những cụm từ  dài, những mệnh đề; không viết tắt những cụm từ ít xuất hiện trong báo cáo.
Nếu cần viết tắt những từ thuật ngữ, tên các cơ quan, tổ chức,\ldots thì được viết tắt sau lần viết thứ nhất có kèm theo chữ viết tắt trong ngoặc đơn.
Nếu báo cáo có nhiều chữ viết tắt thì phải có bảng danh mục các chữ viết tắt (xếp theo thứ tự ABC) ở phần đầu báo cáo.

Nhắc lại: \textbf{không lạm dụng việc viết tắt} trong báo cáo.
Khi các bạn sử dụng từ viết tắt, người đọc sẽ phải lật lại những phần đã đọc, để tìm lại xem từ viết tắt đó nghĩa là gì.
Việc này sẽ làm chậm tốc độ đọc và sẽ khiến người đọc khó theo dõi báo cáo của bạn hơn.
Nếu có thể, hạn chế hoàn toàn việc dùng viết tắt.

\section{Tài liệu tham khảo và cách trích dẫn}

Mọi ý kiến, khái niệm có ý nghĩa, mang tính chất gợi ý không phải của riêng tác giả và mọi tham khảo khác phải được trích dẫn và chỉ ra nguồn trong danh mục Tài liệu tham khảo của báo cáo. Nguồn được trích dẫn phải được liệt kê chính xác trong danh mục Tài liệu tham khảo.

Việc trích dẫn, tham khảo chủ yếu nhằm thừa nhận nguồn của những ý tưởng có giá trị giúp người đọc theo được mạch suy nghĩ của tác giả, không làm trở ngại việc đọc.


Không trích dẫn những kiến thức phổ biến, mọi người đều biết cũng như không làm báo cáo nặng nề với những tham khảo trích dẫn.

Nếu không có điều kiện tiếp cận được một tài liệu gốc mà phải trích dẫn thông qua một tài liệu khác thì phải nêu ra trích dẫn này, đồng thời tài liệu gốc đó không được liệt kê trong danh mục tài liệu tham khảo của báo cáo.

\chapter{Hướng dẫn sử dụng template}
\label{Chapter3}

\section{Trích dẫn tài liệu}

Dùng lệnh $\textbackslash cite$ để trích dẫn một hoặc nhiều tài liệu tham khảo.
Tài liệu tham khảo có thể là trang web~\cite{Listings,HDLVThS}, bài báo khoa học~\cite{1994-Cavnar}, sách~\cite{1984-TeX-Knuth,2006-DDien,2006-NPTV}, bài tạp chí~\cite{1989-TED} hoặc các nguồn tham khảo khác. 
Lưu ý khi trích dẫn tài liệu tham khảo, cần viết câu sao cho bỏ phần trong cặp ngoặc vuông đi thì câu vẫn đầy đủ ý nghĩa.
Ví dụ, thay vì viết ``Nghiên cứu~\cite{1989-TED} chỉ ra rằng ... '' thì nên viết ``Nghiên cứu của Zhang~\cite{1989-TED} chỉ ra rằng ...''.
Một ví dụ khác, thay vì viết ``... như trong công trình nghiên cứu~\cite{1994-Cavnar}.'' thì nên viết ``... như trong công trình nghiên cứu của Cavnar và Trenkle~\cite{1994-Cavnar}.''

\section{Chèn mã nguồn}

Để chèn mã nguồn, cần dùng package listings~\cite{Listings}:

\begin{lstlisting}
\usepackage{listings}
\end{lstlisting}

Mã nguồn có thể được chèn trực tiếp như sau:

\begin{lstlisting}[language=Python]
print "Hello, World!"
\end{lstlisting}

hoặc chèn thông qua tập tin chứa mã nguồn trong thư mục $SourceCode$ như sau:

\lstinputlisting[language=C++]{SourceCode/hello.cpp}

Để chèn mã giả, cần dùng package algorithm~\cite{Algorithm}:

\begin{lstlisting}
\usepackage{algorithm}
\end{lstlisting}

Có thể chèn mã giả vào như sau:

\begin{algorithm}
\caption{My algorithm}\label{euclid}
\begin{algorithmic}[1]
\Procedure{MyProcedure}{}
\State $\textit{stringlen} \gets \text{length of }\textit{string}$
\State $i \gets \textit{patlen}$
\BState \emph{top}:
\If {$i > \textit{stringlen}$} \Return false
\EndIf
\State $j \gets \textit{patlen}$
\BState \emph{loop}:
\If {$\textit{string}(i) = \textit{path}(j)$}
\State $j \gets j-1$.
\State $i \gets i-1$.
\State \textbf{goto} \emph{loop}.
\State \textbf{close};
\EndIf
\State $i \gets i+\max(\textit{delta}_1(\textit{string}(i)),\textit{delta}_2(j))$.
\State \textbf{goto} \emph{top}.
\EndProcedure
\end{algorithmic}
\end{algorithm}

\section{Hình ảnh}

Để chèn hình ảnh, cần dùng package graphicx~\cite{Figures}:

\begin{lstlisting}
\usepackage{graphicx}
\end{lstlisting}

Hình \ref{fig:vd1}, hình \ref{fig:vd2} là một số ví dụ về chèn hình ảnh.

\begin{figure}[htp]
\centering
\includegraphics[width=6 cm]{images/logo-khtn.png}
\caption{Hình ví dụ 1}
\label{fig:vd1}
\end{figure}

\begin{figure*}[htp]
\centering
\includegraphics[width=40 mm]{images/logo-khtn.png}
\caption{Hình ví dụ 2}
\label{fig:vd2}
\end{figure*}

\section{Bảng biểu}

Để tạo bảng biểu, tham khảo thêm tại sharelatex.com~\cite{Tables}.
Bảng~\ref{tab:vd1} là một ví dụ về bảng.
Ngoài ra, có một số tool online~\footnote{Như trang http://www.tablesgenerator.com/} có thể được dùng để tạo bảng biểu một cách trực quan.

\begin{table}[ht]
\caption{Bảng ví dụ 1}

\label{tab:vd1}%
\begin{center}
\begin{tabular}{ |p{3cm}||p{3cm}|p{3cm}|p{3cm}|  }
 \hline
 \multicolumn{4}{|c|}{Country List} \\
 \hline
 Country Name     or Area Name& ISO ALPHA 2 Code &ISO ALPHA 3 Code&ISO numeric Code\\
 \hline
 Afghanistan   & AF    &AFG&   004\\
 Aland Islands&   AX  & ALA   &248\\
 Albania &AL & ALB&  008\\
 Algeria    &DZ & DZA&  012\\
 American Samoa&   AS  & ASM&016\\
 Andorra& AD  & AND   &020\\
 Angola& AO  & AGO&024\\
 \hline
\end{tabular}
\end{center}
\end{table}

\section{Công thức}

Công thức có thể chèn vào trong cùng một dòng như $ \sqrt{a^2+b^2} $ hoặc nằm trên dòng riêng như công thức~\ref{eu_eqn}.

\begin{equation} \label{eu_eqn}
x = a_0 + \frac{1}{a_1 + \frac{1}{a_2 + \frac{1}{a_3 + a_4}}}
\end{equation}

% Công trình của tác giả (nếu không có thì comment 02 dòng dưới)
\addcontentsline{toc}{chapter}{Danh mục công trình của tác giả}
\chapter*{Danh mục công trình của tác giả}
\label{Appendix1}

\begin{enumerate}
\item Tạp chí ABC
\item Tạp chí XYZ
\end{enumerate}

% In tài liệu tham khảo
\addcontentsline{toc}{chapter}{Tài liệu tham khảo}
\printbibheading[title={Tài liệu tham khảo}]

\printbibliography[heading=subbibliography, title={Tiếng Việt}, keyword=Viet, resetnumbers=true]

\DeclareNameAlias{sortname}{last-first}
\DeclareNameAlias{default}{last-first}

\printbibliography[heading=subbibliography, title={Tiếng Anh}, notkeyword=Viet, resetnumbers=4] 
% ===================================================================== %
% CHÚ Ý: phải gán lại resetnumbers=số tài liệu tham khảo tiếng Việt + 1 %
% ===================================================================== %

% Phần phụ lục
\appendix

\chapter{Tìm kiếm siêu tham số}
\label{Appendix1}

Số lượng siêu tham số trong một hệ thống FL là quá lớn. Do đó, sau khi giới hạn không gian tìm kiếm thì số phép thử còn lại cũng khó có thể thực hiện vét cạn. Hệ thống FL trình bày trong khóa luận cũng không nằm ngoại lệ.

Các siêu tham số của hệ thống FL của khóa luận bao gồm: số máy khách tham gia huấn luyện trong một bước huấn luyện toàn cục (\#clients/round), số bước huấn luyện cục bộ (\#epochs), số bước huấn luyện toàn cục (\#rounds), lượng dữ liệu trong một batch dữ liệu (batch_size), số lớp phần riêng đối với các thuật toán sử dụng kỹ thuật PL (\#per_layer) và các siêu tham số học sử dụng trong tối ưu mạng học sâu bằng kỹ thuật SGD.

Để phù hợp cho phần cứng của các máy khách có cấu hình yếu trong kịch bản Horizontal FL, số bước huấn luyện cục bộ và lượng dữ liệu trong một batch dữ liệu được giữ cố định lần lượt là 1 và 32. Từ việc khảo sát các thí nghiệm FL của các nghiên cứu gần đây, số máy khách tham gia huấn luyện toàn cục được được chọn lần lượt là 2, 5 và 10 máy. Trong đó, sử dụng 5 máy khách tham gia huấn luyện cùng lúc cho kết quả cao hơn một chút so với việc sử dụng 2 hay 10 máy và tiêu tốn chi phí tính toán ở một mức chấp nhận được. Do giới hạn phần cứng và theo quan sát quá trình hội tụ, số lượng bước huấn luyện toàn cục được giữ ở mức 300 cho tập MNIST và 600 cho tập CIFAR-10.

Đối với kích thước mạng học sâu được cài đặt trong khóa luận, việc duy trì một lớp phần chung và một lớp phần riêng cho mạng neural MNIST là tất nhiên. Với mạng học sâu dùng cho tập CIFAR-10, số lớp phần riêng được chọn lần lượt là 1, 2 và 3 lớp (tính từ lớp tuyến tính cuối cùng). Kết quả chạy thực nghiệm có thấy, việc sử dụng lớp tuyến tính cuối cùng làm phần riêng và các lớp học sâu còn lại làm phần chung cho kết quả tốt nhất trên tập CIFAR-10.

Ngoại trừ siêu tham số học của từng thuật toán, các siêu tham số kể trên đều được giữ cố định trong quá trình huấn luyện. Bảng \ref{tab:fixed_hyper_param} trình bày tóm tắt các giá trị siêu tham số này.

\begin{table}[H]
    \centering
    \caption{Bảng các siêu tham số cố định của hệ thống trên MNIST và CIFAR-10}
    \label{tab:fixed_hyper_param}
    \resizebox{\linewidth}{!}{%
    \begin{tabular}{l|ccccc} 
    \toprule
             & \#clients/round    & \#epochs           & \#rounds & batch\_size         & \multicolumn{1}{l}{\#per\_layer}  \\ 
    \hline
    MNIST    & \multirow{2}{*}{5} & \multirow{2}{*}{1} & 300      & \multirow{2}{*}{32} & \multirow{2}{*}{1}                \\
    CIFAR-10 &                    &                    & 600      &                     &                                   \\
    \bottomrule
    \end{tabular}
    }
\end{table}

Các siêu tham số học được tìm kiếm trong khoảng $(10^{-5}, 0.01)$ cho từng thuật toán. Kết quả tìm kiếm được trình bày trong Bảng \ref{tab:hyper_param}.

\begin{table}[H]
    \centering
    \caption{Bảng siêu tham số được sử dụng cho từng thuật toán}
    \label{tab:hyper_param}
    \resizebox{\linewidth}{!}{%
    \begin{tabular}{l|cc} 
    \toprule
    \begin{tabular}[c]{@{}l@{}}\\\end{tabular} & CIFAR-10           & MNIST                                             \\ 
    \hline
    FedAvg, FedAvgMeta                         & $10^{-4}$          & $10^{-5}$                                      \\
    FedPer, FedPerMeta                         & $10^{-4}$          & $10^{-5}$                                      \\
    FedMeta(MAML) ($\alpha,\beta$)             & $(0.01, 0.001)$    & $(0.001, 0.001)$                                \\
    FedMeta(Meta-SGD)($\alpha,\beta$)          & $(0.001, 0.001)$   & $(0.001, 5\times 10^{-4})$  \\
    FedMeta-Per(MAML)($\alpha,\beta$)          & $(0.001, 0.005)$   & $(0.001, 0.001)$                                \\
    FedMeta-Per(Meta-SGD)($\alpha,\beta$)      & $(0.01,0.01)$      & $(0.001, 5\times 10^{-4})$  \\
    \bottomrule
    \end{tabular}
    }
\end{table}
\chapter{Ngữ pháp tiếng Nôm}
\label{Appendix2}

Đây là phụ lục 2.

\end{document} 